\section{The Wilson--Loewner realization hypothesis}\label{sec:realization}

\subsection{An observable before a driver}

A physical candidate must specify boundary input and output spaces,
a field theory and regulator, endpoint states, a reference color
sector, a physical adjoint, a family or ensemble of contours, and a
normalization. After a fixed identification of its boundary spaces
with $L^2(I_L)$, denote its averaged transfer by
$\mathcal V_{\omega,L}$. It must act linearly on arbitrary boundary
inputs, not merely return one scalar expectation for a single curve.

\begin{hypothesis}[Prescribed effective evolution]\label{hyp:realization}
There is an independently specified Wilson boundary construction,
with contours generated or parametrized by a Loewner evolution,
whose normalized transfer is strongly initialized by
$\mathcal V_{0,L}=I$ and satisfies
\begin{equation}\label{eq:physical-flow}
 \partial_\omega\mathcal V_{\omega,L}
       =-G_{\omega,L}\mathcal V_{\omega,L}
\end{equation}
on the arithmetic smooth core, in a causal evolution class for
which this initial-value problem is unique.
\end{hypothesis}

Under these conditions $\mathcal V_{\omega,L}=V_{\omega,L}$.
The existence of the construction, closure into the specified
$G_{\omega,L}$, and the applicable uniqueness statement are all
obligations of the hypothesis. Declaring $-G$ to be an abstract
connection gives a formal ordered transport, but supplies none of
the field or endpoint data. At fixed $L$ the causal convolution
generators formally commute on a common admissible core; the nontrivial question remains how
a color-ordered field observable reduces to this boundary evolution.

The differential equation is intended to \emph{select} admissible
physical and geometric data. There may be no suitable scalar
driver within a proposed family, or several such families. An
explicit real driver or a driver law would be an output of solving
that matching problem. We do not assert that the arithmetic
equation already defines one.

\subsection{The insertion identity that must close}

For smooth matrix fields on a regular curve, set
\[
 \mathcal M=\ii A_\mu\dot X^\mu+\Phi v,\qquad
 \Phi=n^I X_I,\qquad v=|\dot X|.
\]
Ordered Wilson transport obeys
$\partial_sU(s,0)=\mathcal M(s)U(s,0)$.
For a shape change $\eta=\delta X$ and a change $\delta n$,
its exact variation is
\begin{align}\label{eq:variation-summary}
 \delta U(1,0)
 &=C_1U(1,0)-U(1,0)C_0
   +\int_0^1U(1,s)\mathcal I(s)U(s,0)\dd s,\\
 \mathcal I
 &=\ii F_{\nu\mu}\eta^\nu\dot X^\mu
   +vD_\nu\Phi\,\eta^\nu
   +\Phi\frac{\dot X\cdot\dot\eta}{v}
   +vX_I\delta n^I,\qquad C_s=\ii A_\nu\eta^\nu.
 \notag
\end{align}
The boundary terms combine with the variations of actual endpoint
fields and their polarizations. The proof, conventions, tangent-map
version and endpoint formula are in \cite[Section~S4]{Supplement}.

To obtain \eqref{eq:physical-flow}, the averaged insertions and all
normalization terms must act on boundary inputs as
$-G_{\omega,L}\mathcal V_{\omega,L}$. Field equations, Ward
identities or a controlled reduction must establish this closure.
A universal formula depending only on the curve cannot do so:
in the abelian example $A_x=cy$, a straight line at $y=0$ has
$U=1$ for every $c$, but its transverse variation depends on $c$.
This obstruction concerns bare transport; it leaves open a
field-dependent averaged closure.

\subsection{Three parameters that must be related explicitly}

Use the chordal equation
\begin{equation}\label{eq:main-loewner}
 \partial_tg_t(z)=\frac2{g_t(z)-u_t},
 \qquad \operatorname{hcap}=2t,
\end{equation}
where $u_t$ is a real driver, distinguished from Wilson transport
$U$. For a smooth reference curve $\zeta(s)$ in a region where
$f_t=g_t^{-1}$ is regular,
\[
 X_t(s)=f_t(\zeta(s)),\qquad
 \partial_tX_t(s)=-\frac{2f_t'(\zeta(s))}{\zeta(s)-u_t}.
\]
These velocities can be inserted into
\eqref{eq:variation-summary}; they do not define transport at a
singular growing tip. A stochastic trace would need additional
regularization and a defined limiting observable.

At fixed $L$, a dictionary $\omega=\omega(t)$ would require
\begin{equation}\label{eq:clock-dictionary}
 \partial_t\mathcal V_{\omega(t),L}
 =-\dot\omega(t)G_{\omega(t),L}\mathcal V_{\omega(t),L}.
\end{equation}
Loewner capacity $t$, arithmetic shift $\omega$, and support length
$L$ have different meanings. Even the explicit constant-driver
semicircle dictionary in \cite[Section~S3]{Supplement} identifies
neither $t$ nor the logarithmic endpoint radius with $\omega$.
When $L$ changes, so do the space and compression; the fixed-window
equation \eqref{eq:clock-dictionary} omits that change.
Section~\ref{sec:path} treats it as a separate finite gluing step.
